\documentclass[10pt,a4paper]{amsart}
\usepackage[margin=19mm]{geometry}
\usepackage{amsmath,amssymb,mathtools}
\usepackage[scr=rsfs,cal=euler]{mathalfa}
\usepackage{xfrac}
\usepackage[unicode,hidelinks,bookmarks=false]{hyperref}
\newcommand{\ie}{i.e.\ }
\newcommand{\cf}{cf.\ }
\newcommand{\resp}{respectively\ }
\newcommand{\Subsection}{\subsection}
\providecommand{\nobreakdash}{\nobreak\hbox{-}}
\setlength{\emergencystretch}{2em}
\multlinegap0pt
\usepackage{bm}
\usepackage{bbm}
\usepackage{dsfont}
\usepackage{xspace}
\usepackage{cancel}
\usepackage{mathtools}
\usepackage[shortlabels]{enumitem}
\usepackage{amsthm}
\makeatletter
\@ifundefined{thrm}{\newtheorem{thrm}{Thrm}}{}
\@ifundefined{lem}{\newtheorem{lem}[thrm]{Lemma}}{}
\makeatother
\newcommand{\A}{\mathbb{A}}
\newcommand{\N}{\mathbb{N}}
\newcommand{\Q}{\mathbb{Q}}
\newcommand{\Z}{\mathbb{Z}}
\newcommand{\ba}{\boldsymbol{a}}
\newcommand{\bb}{\boldsymbol{b}}
\newcommand{\bx}{\boldsymbol{x}}
\newcommand{\frm}{\mathfrak{m}}
\newcommand{\mA}{\mathcal{A}}
\newcommand{\mI}{\mathcal{I}}
\newcommand{\mJ}{\mathcal{J}}
\newcommand{\mV}{\mathcal{V}}
\newcommand{\oA}{\overline{A}}
\newcommand{\oX}{\mkern 1.5mu\overline{\mkern-1.5mu X \mkern-1.5mu}\mkern 1.5mu}
\newcommand{\ox}{\overline{x}}
\newcommand{\tc}{\widetilde{c}}
\title[S-unit equations in modules and linear-exponential Diophanti]{S-unit equations in modules and linear-exponential Diophantine equations}
\author{Ruiwen Dong, Doron Shafrir}
\date{Source: arXiv:2505.19141v2 (CC BY 4.0)}
\begin{document}
\maketitle

\noindent\emph{Source and licence.} S-unit equations in modules and linear-exponential Diophantine equations. Version arXiv:2505.19141v2 (CC BY 4.0), distributed under the Creative Commons Attribution 4.0 International licence, as recorded in the arXiv licence metadata for this exact version and shown on the abstract page https://arxiv.org/abs/2505.19141v2. Full rights notice: licenses/arxiv-2505.19141-CC-BY-4.0.txt.\par\smallskip

\noindent\textit{Editorial scope.} Counted unit: the Lem whose source label is \texttt{lem:syminduction} in the article (source0.tex lines 3689--3692, source label \texttt{lem:syminduction}), delivered together with the article's own complete original proof of it (source0.tex lines 3693--3888, one \texttt{proof} environment of 196 lines). No mathematics is written, repaired or re-derived: statement and proof are the article's own text line for line. Required context retained uncounted: lem:lifting (lines 1156--1159); eq:defbxgeq (lines 3580--3582); lem:lindepmodp (lines 3589--3601); lem:powpn (lines 3663--3668). Every definition, display and label of those lines is retained unchanged. Omitted from this package and not redistributed: 1--1155, 1160--3579, 3583--3588, 3602--3662, 3669--3688, 3889--4359. Nothing inside a retained excerpt is omitted or altered: every retained body is delivered exactly as the article's lines are written, so no omission receipt of any kind accompanies this package. Citations: the retained excerpts carry no citation command at all, so no bibliography entry of the article is transcribed and no reference is redistributed. Reference adaptation: none was needed and none was made; every reference inside the delivered unit resolves inside it, so no unresolved reference or citation is produced. Printed slips kept literal: no printed word, symbol, hypothesis or number of the counted statement or of its proof was altered. Layout and class: the article's own class is \texttt{article} with packages including geometry, footmisc, enumitem, appendix, xcolor, array, amsthm, amsmath, amssymb, amsfonts, graphicx, thm-restate, hyperref, mathtools; the wrapper is the lane's pinned \texttt{amsart} editorial wrapper at 10pt on A4, which restates the theorem-like environments the retained text needs and adds the article's own macro definitions that the retained text needs (\texttt{\textbackslash A}, \texttt{\textbackslash N}, \texttt{\textbackslash Q}, \texttt{\textbackslash Z}, \texttt{\textbackslash ba}, \texttt{\textbackslash bb}, \texttt{\textbackslash bx}, \texttt{\textbackslash frm}, \texttt{\textbackslash mA}, \texttt{\textbackslash mI}, \texttt{\textbackslash mJ}, \texttt{\textbackslash mV}, \texttt{\textbackslash oA}, \texttt{\textbackslash oX}, \texttt{\textbackslash ox}, \texttt{\textbackslash tc}), with extra wrapper packages (bm, bbm, dsfont, xspace, cancel, mathtools, enumitem). The delivered numbering is the unit's own. Assets: no retained span contains a figure environment, a graphics-inclusion command, a PDF-page inclusion, a table or a TikZ picture; no image file is copied into this package, the package declares no asset dependency and no image file is redistributed with it. Licence: version arXiv:2505.19141v2 of this article is distributed under the Creative Commons Attribution 4.0 International licence, as recorded in the arXiv licence metadata for this exact version and shown on the abstract page https://arxiv.org/abs/2505.19141v2. A full rights notice is delivered at licenses/arxiv-2505.19141-CC-BY-4.0.txt. Characters: every retained excerpt is pure ASCII, so the delivered engine, XeLaTeX with its default Latin Modern text font, reports no missing character.
\begin{lem}\label{lem:lifting}
    Let $R$ be a ring, and $P$ be an ideal such that $p \in P$, and $P^t = 0$ for some $t \in \N$.
    Let $f, g \in R$ such that $f - g \in P$. Then $f^{p^{r}} = g^{p^{r}}$ for all $r \geq t-1$.
\end{lem}

\begin{equation}\label{eq:defbxgeq}
\bx_i^{(\geq r)} \coloneqq (\ox^{p^{\ell r} \ba_i + \bb_i}, \ox^{p^{\ell (r+1)} \ba_i + \bb_i}, \ox^{p^{\ell (r+2)} \ba_i + \bb_i}, \ldots) \in \A^{\N}.
\end{equation}

\begin{lem}\label{lem:lindepmodp}
    Let $\A$ be a field of characteristic $p$.
    Let $\ox = (x_1, \ldots, x_N)$ be a tuple of non-zero elements in $\A$.
    Pick a maximal subset $\mI \subseteq \{1, \ldots, K\}$ such that $\bx_i^{(\geq r)}, i \in \mI$ are $\A$-linearly independent for all $r$.
    Then for any $j \notin \mI$, we have
    \begin{equation}\label{eq:linconda}
        \ox^{\ba_j + \bb_j} = \sum_{i \in \mI} c_{j,i} \cdot \ox^{\ba_i + \bb_i},
    \end{equation}
    for some $c_{j,i} \in \A, i \in \mI$ satisfying
    \begin{equation}\label{eq:lincondc}
        c_{j,i}^{p^{\ell}} \cdot \ox^{(p^{\ell} - 1) \bb_i} = c_{j,i} \cdot \ox^{(p^{\ell}-1)\bb_j}.
    \end{equation}
\end{lem}

\begin{lem}\label{lem:powpn}
    Let $f \in \Z[\oX^{\pm}]$, then for all $n \geq t$, we have 
    \[
    f^{p^{\ell n}}(Y_1, \ldots, Y_k) \equiv f^{p^{\ell t}}(Y_1^{p^{\ell (n-t)}}, \ldots, Y_k^{p^{\ell (n-t)}}) \mod p^t.
    \]
\end{lem}

\begin{lem}\label{lem:syminduction}
    Let $s \leq t$ be an integer, and let $v_1, \ldots, v_K \in \frm^{s}\mV$. 
    Suppose there exists $m \in \N$ such that $\sum_{i = 1}^K \oA^{p^{\ell n} \ba_i + \bb_i} v_i = 0$ holds for all $n \geq m$. Then $\sum_{i = 1}^K \oA^{p^{\ell n} \ba_i + \bb_i} v_i = 0$ holds for all $n \geq (t+1-s)t$.
\end{lem}

\begin{proof}
    We use reverse induction on $s$, starting from $t$ and gradually decreasing to $0$.
    The case where $s = t$ is trivial because $\frm^{t}\mV = 0$. We now focus on the induction step.
    Suppose the statement is true for $s$, we show it for $s-1$.
    Let $v_1, \ldots, v_K \in \frm^{s-1}\mV$.
    
    If $v_i \in \frm^{s} \mV$ for all $i = 1, \ldots, K$, then we conclude directly using the induction hypothesis for $s$.
    Suppose now that $v_i \notin \frm^{s} \mV$ for some $i$. 
    Without loss of generality, we can suppose $v_1 \in \frm^{s-1} \mV \setminus \frm^{s} \mV, \ldots, v_{K'} \notin \frm^{s-1} \mV \setminus \frm^{s} \mV$, and $v_{K'+1} \in \frm^{s} \mV, \ldots, v_K \in \frm^{s} \mV$, where $1 \leq K' \leq K$.
    %By a slight abuse of notation, for $i = K'+1, \ldots, K$, let $\frac{v_i}{p}$ denote the element in $\mV$ such that $p \cdot \frac{v_i}{p} = \mV$.
    
    Let $\A \coloneqq \mA/\frm$, it is a field of characteristic $p$.
    Let
    \[
    x_1 \coloneqq A_1 + \frm, x_2 \coloneqq A_2 + \frm, \ldots, x_N \coloneqq A_N + \frm,
    \]
    be elements of $\A$. These are non-zero since $A_1, \ldots, A_N$ are invertible in $\mA$. Denote $\ox \coloneqq (x_1, \ldots, x_N)$.
    
    As in Equation~\eqref{eq:defbxgeq}, for each $i = 1, \ldots, K$, and $r \in \N$, denote the sequence
    \[
    \bx_i^{(\geq r)} \coloneqq (\ox^{p^{\ell r} \ba_i + \bb_i}, \ox^{p^{\ell (r+1)} \ba_i + \bb_i}, \ox^{p^{\ell (r+2)} \ba_i + \bb_i}, \ldots) \in \A^{\N}.
    \]
    Let $\mI \subseteq \{1, \ldots, K'\}$ be a maximal subset such that $\bx_i^{(\geq r)}, i \in \mI$ are $\A$-linearly independent for all $r \in \N$.
    Without loss of generality suppose $\mI = \{1, \ldots, k\}$ for some $k \leq K'$.
    By Lemma~\ref{lem:lindepmodp}, we can write
    \begin{multline*}
    \ox^{\ba_{k+1} + \bb_{k+1}} = c_{k+1,1} \cdot \ox^{\ba_1 + \bb_1} + \cdots + c_{k+1,k} \cdot \ox^{\ba_k + \bb_k}, \\
    \ddots \\
    \ox^{\ba_{K'} + \bb_{K'}} = c_{K',1} \cdot \ox^{\ba_1 + \bb_1} + \cdots + c_{K',k} \cdot \ox^{\ba_k + \bb_k},
    \end{multline*}
    for some $c_{k+1, 1}, \ldots, c_{K', k} \in \A$ that satisfy
    \begin{equation}\label{eq:condc}
    c_{j,i}^{p^{\ell}} \cdot \ox^{(p^{\ell} - 1) \bb_i} = c_{j,i} \cdot \ox^{(p^{\ell}-1)\bb_j}.
    \end{equation}
    For $j = k+1, \ldots, K'; i = 1, \ldots, k$, take any $\tc_{j, i} \in \mA$ such that $\tc_{j, i} = c_{j, i} + \frm$.
    This yields
    \begin{multline*}
    \oA^{\ba_{k+1} + \bb_{k+1}} \equiv \tc_{k+1,1} \oA^{\ba_1 + \bb_1} + \cdots + \tc_{k+1,k} \oA^{\ba_k + \bb_k} \mod \frm, \\\;\ddots\; \\
    \ox^{\ba_{K'} + \bb_{K'}} \equiv \tc_{K',1} \oA^{\ba_1 + \bb_1} + \cdots + \tc_{K',k} \oA^{\ba_k + \bb_k} \mod \frm.
    \end{multline*}
    
    Applying Lemma~\ref{lem:lifting} to the ring $\mA$ and the ideal $\frm$, we have
    \begin{multline*}
    \oA^{p^{\ell n}(\ba_{k+1} + \bb_{k+1})} = \left(\tc_{k+1,1} \oA^{\ba_1 + \bb_1} + \cdots + \tc_{k+1,k} \oA^{\ba_k + \bb_k} \right)^{p^{\ell n}}, \\\;\ddots\; \\
    \oA^{p^{\ell n}(\ba_{K'} + \bb_{K'})} = \left(\tc_{K',1} \oA^{\ba_1 + \bb_1} + \cdots + \tc_{K',k} \oA^{\ba_k + \bb_k} \right)^{p^{\ell n}},
    \end{multline*}
    for all $\ell n \geq t-1$.
    Then for all $n \geq t-1$, the equation $\sum_{i = 1}^K \oA^{p^{\ell n} \ba_i + \bb_i} v_i = 0$ is equivalent to
    \begin{multline}\label{eq:sumyv}
    \sum_{i = 1}^k \oA^{p^{\ell n} \ba_i + \bb_i} v_i + \sum_{j = k+1}^{K'} \left(\tc_{j,1} \oA^{\ba_1 + \bb_1} + \cdots + \tc_{j,k} \oA^{\ba_k + \bb_k} \right)^{p^{\ell n}} \oA^{(1 - p^{\ell n}) \bb_j} v_j + \sum_{i = K'+1}^K \oA^{p^{\ell n} \ba_i + \bb_i} v_i = 0.
    \end{multline}
    For $j = k+1, \ldots, K'$, consider the polynomial $f(Y_1, \ldots, Y_k) \coloneqq Y_1 + \cdots + Y_k \in \Z[Y_1, \ldots, Y_k]$.
    By Lemma~\ref{lem:powpn} we have
    $
    f^{p^{\ell n}}(Y_1, \ldots, Y_k) \equiv f^{p^{\ell t}}(Y_1^{p^{\ell (n-t)}}, \ldots, Y_k^{p^{\ell (n-t)}}) \mod p^t,
    $
    for all $n \geq t$.
    We can write the polynomial $f^{p^{\ell t}}(Y_1, \ldots, Y_k) = (Y_1 + \cdots + Y_k)^{p^{\ell t}}$ in the form
    \[
    Y_1^{p^{\ell t}} + \cdots + Y_k^{p^{\ell t}} + p \sum_{i \in \mJ} u_{i} Y_1^{d_{i1}} Y_2^{d_{i2}} \cdots Y_k^{d_{ik}}
    \]
    for some finite index set $\mJ$, and with $u_i \in \Z, d_{i1} + \cdots + d_{ik} = p^{\ell t},$ for all $i \in \mJ$.
    Then
    \begin{equation}\label{eq:cY}
    (Y_1 + \cdots + Y_k)^{p^{\ell n}} \equiv Y_1^{p^{\ell n}} + \cdots + Y_k^{p^{\ell n}} + p \sum_{i \in \mJ} u_{i} Y_1^{p^{\ell (n - t)} d_{i1}} Y_2^{p^{\ell (n - t)} d_{i2}} \cdots Y_k^{p^{\ell (n - t)} d_{ik}} \mod p^t.
    \end{equation}
    For each $j = k+1, \ldots, K'$, consider the ring homomorphism from $\Z[Y_1, \ldots, Y_k]$ to $\mA$, that sends $Y_1, Y_2, \ldots, Y_k$ respectively to $\tc_{j,1} \oA^{\ba_1 + \bb_1}, \tc_{j,2} \oA^{\ba_2 + \bb_2}, \ldots, \tc_{j,k} \oA^{\ba_k + \bb_k}$.
    Since $p^t = 0$ in $\mA$, applying this homomorphism to Equation~\eqref{eq:cY} yields
    \begin{multline}\label{eq:unfoldpln}
    \left(\tc_{j,1} \oA^{\ba_1 + \bb_1} + \cdots + \tc_{j,k} \oA^{\ba_k + \bb_k}\right)^{p^{\ell n}} = \tc_{j,1}^{p^{\ell n}} \cdot \oA^{p^{\ell n}(\ba_1 + \bb_1)} + \cdots + \tc_{j,k}^{p^{\ell n}} \cdot \oA^{p^{\ell n}(\ba_k + \bb_k)} + \\
    p \sum_{i \in \mJ} u_{i} \tc_{j,1}^{p^{\ell (n - t)} d_{i1}} \cdot \oA^{p^{\ell (n-t)}d_{i1}(\ba_1 + \bb_1)} \cdots \tc_{j,k}^{p^{\ell (n - t)} d_{ik}} \cdot \oA^{p^{\ell (n-t)}d_{ik}(\ba_k + \bb_k)}.
    \end{multline}
    
    Recall from Equation~\eqref{eq:condc} that $c_{j,i}^{p^{\ell}} \cdot \ox^{(p^{\ell} - 1) \bb_i} = c_{j,i} \cdot \ox^{(p^{\ell}-1)\bb_j}$ for all $j, i$.
    So
    \[
    \tc_{j,i}^{p^{\ell}} \cdot \oA^{(p^{\ell} - 1) \bb_i} \equiv \tc_{j,i} \cdot \oA^{(p^{\ell}-1)\bb_j} \mod \frm.
    \]
    For any $n \geq t$, taking $p^{\ell (n-1)}$-th power, $p^{\ell (n-2)}$-th power, $\ldots,$ $p^{\ell (t-1)}$-th power, on both sides and using Lemma~\ref{lem:lifting} yield respectively
    \begin{align*}
    \tc_{j,i}^{p^{\ell n}} \cdot \oA^{(p^{\ell n} - p^{\ell (n-1)}) \bb_i} & = \tc_{j,i}^{p^{\ell (n-1)}} \cdot \oA^{(p^{\ell n} - p^{\ell (n-1)})\bb_j}, \\
    \tc_{j,i}^{p^{\ell (n-1)}} \cdot \oA^{(p^{\ell (n-1)} - p^{\ell (n-2)}) \bb_i} & = \tc_{j,i}^{p^{\ell (n-2)}} \cdot \oA^{(p^{\ell (n-1)} - p^{\ell (n-2)})\bb_j}, \\
    & \vdots \\
    \tc_{j,i}^{p^{\ell t}} \cdot \oA^{(p^{\ell t} - p^{\ell (t-1)}) \bb_i} & = \tc_{j,i}^{p^{\ell (t-1)}} \cdot \oA^{(p^{\ell t} - p^{\ell (t-1)})\bb_j}.
    \end{align*}
    Their product yields for all $n \geq t$,
    \begin{equation}\label{eq:tc1}
    \tc_{j,i}^{p^{\ell n}} = \tc_{j,i}^{p^{\ell (t-1)}} \cdot \oA^{(p^{\ell n} - p^{\ell (t-1)})(\bb_j - \bb_i)}.
    \end{equation}
    For all $n \geq 2t$, replacing $n$ with $n-t$ in Equation~\eqref{eq:tc1} yields
    \begin{equation}\label{eq:tc2}
    \tc_{j,i}^{p^{\ell (n-t)}} = \tc_{j,i}^{p^{\ell (t-1)}} \cdot \oA^{(p^{\ell (n-t)} - p^{\ell (t-1)})(\bb_j - \bb_i)}.
    \end{equation}
    Substituting the terms $\tc_{j,i}^{p^{\ell n}}$ and $\tc_{j,i}^{p^{\ell (n-t)}}$ in the right hand side of Equation~\eqref{eq:unfoldpln} using Equations~\eqref{eq:tc1} and \eqref{eq:tc2}, we obtain
    \begin{multline}\label{eq:sub1}
        \left(\tc_{j,1} \oA^{\ba_1 + \bb_1} + \cdots + \tc_{j,k} \oA^{\ba_k + \bb_k}\right)^{p^{\ell n}} = \\
        \left(\tc_{j,1}^{p^{\ell (t-1)}} \cdot \oA^{p^{\ell n}\ba_1 + p^{\ell (t-1)} \bb_1} + \cdots + \tc_{j,k}^{p^{\ell (t-1)}} \cdot \oA^{p^{\ell n}\ba_k + p^{\ell (t-1)} \bb_k} \right) \cdot \oA^{(p^{\ell n} - p^{\ell (t-1)}) \bb_j} + \\
        p \sum_{i \in \mJ} \left( u_{i} \tc_{j,1}^{p^{\ell (t-1)}d_{i1}} \oA^{\left(p^{\ell (n-t)}\ba_1 + p^{\ell (t-1)} \bb_1  \right)d_{i1}} \cdots \tc_{j,k}^{p^{\ell (t-1)} d_{ik}} \oA^{\left(p^{\ell (n-t)}\ba_k + p^{\ell (t-1)} \bb_k \right) d_{ik} } \right) \oA^{(p^{\ell (n-t)} - p^{\ell (t-1)}) \bb_j \cdot p^{\ell t}},
    \end{multline}
    for all $n \geq 2t$.
    Note that the term $p^{\ell t}$ on the final exponent comes from using $d_{i1} + \cdots + d_{ik} = p^{\ell t}$ in the above substitution.
    Using Equation~\eqref{eq:sub1} to substitute $\left(\tc_{j,1} \oA^{\ba_1 + \bb_1} + \cdots + \tc_{j,k} \oA^{\ba_k + \bb_k}\right)^{p^{\ell n}}$ in Equation~\eqref{eq:sumyv} yields
    \begin{multline}\label{eq:sumyv2}
        \sum_{i = 1}^k \oA^{p^{\ell n} \ba_i + \bb_i} v_i \\
        + \sum_{j = k+1}^{K'} \left(\tc_{j,1}^{p^{\ell (t-1)}} \cdot \oA^{p^{\ell n}\ba_1 + p^{\ell (t-1)} \bb_1} + \cdots + \tc_{j,k}^{p^{\ell (t-1)}} \cdot \oA^{p^{\ell n}\ba_k + p^{\ell (t-1)} \bb_k} \right) \cdot \oA^{(1 - p^{\ell (t-1)}) \bb_j} v_j \\
        + p \sum_{j = k+1}^{K'} \sum_{i \in \mJ} u_{i} \tc_{j,1}^{p^{\ell (t-1)}d_{i1}} \cdots \tc_{j,k}^{p^{\ell (t-1)} d_{ik}} \cdot \oA^{p^{\ell (n-t)}(\ba_1 d_{i1} + \cdots + \ba_k d_{ik}) + p^{\ell (t-1)} (\bb_1 d_{i1} + \cdots + \bb_k d_{ik})} \cdot \oA^{(1 - p^{\ell (2t-1)}) \bb_j} v_j \\
        + \sum_{i = K'+1}^K \oA^{p^{\ell n} \ba_i + \bb_i} v_i = 0.
    \end{multline}
    Combining all the terms containing $\oA^{p^{\ell n} \ba_i}$ yields
    \begin{multline}\label{eq:sumyv3}
        \sum_{i = 1}^k \oA^{p^{\ell n} \ba_i + \bb_i} \left(v_i + \sum_{j = k+1}^{K'} \tc_{j,i}^{p^{\ell (t-1)}} \cdot \oA^{(p^{\ell (t-1)} - 1) (\bb_i - \bb_j)} v_j \right) \\
        + p \sum_{j = k+1}^{K'} \sum_{i \in \mJ} u_{i} \tc_{j,1}^{p^{\ell (t-1)}d_{i1}} \cdots \tc_{j,k}^{p^{\ell (t-1)} d_{ik}} \cdot \oA^{p^{\ell (n-t)}(\ba_1 d_{i1} + \cdots + \ba_k d_{ik}) + p^{\ell (t-1)} (\bb_1 d_{i1} + \cdots + \bb_k d_{ik})} \cdot \oA^{(1 - p^{\ell (2t-1)}) \bb_j} v_j \\
        + \sum_{i = K'+1}^K \oA^{p^{\ell n} \ba_i + \bb_i} v_i = 0.
    \end{multline}
    We have shown that the Equation $\sum_{i = 1}^K \oA^{p^{\ell n} \ba_i + \bb_i} v_i = 0$ is equivalent to Equation~\eqref{eq:sumyv3} for $n \geq 2t$.
    
    Note that $p v_{k+1}, \ldots, p v_{K'}, v_{K'+1}, \ldots, v_K \in \frm^s \mV$.
    Therefore for all $n \geq \max\{m, 2t\}$, Equation~\eqref{eq:sumyv3} yields
    \begin{equation}\label{eq:sumyvmod}
    \sum_{i = 1}^k \oA^{p^{\ell n} \ba_i + \bb_i} \left(v_i + \sum_{j = k+1}^{K'} \tc_{j,i}^{p^{\ell (t-1)}} \cdot \oA^{(p^{\ell (t-1)} - 1) (\bb_i - \bb_j)} v_j \right) \in \frm^s \mV.
    \end{equation}
    Note that $v_1, \ldots, v_k \in \frm^{s-1} \mV$, so Equation~\eqref{eq:sumyvmod} is equivalent the following equation in the $\A$-module $\frm^{s-1} \mV/\frm^s\mV$:
    \[
    \sum_{i = 1}^k \ox^{p^{\ell n} \ba_i + \bb_i} \left(v_i + \sum_{j = k+1}^{K'} c_{j,i}^{p^{\ell (t-1)}} \cdot \ox^{(p^{\ell (t-1)} - 1) (\bb_i - \bb_j)} v_j + \frm^s\mV \right) = 0.
    \]

    Note that $\frm^{s-1} \mV/\frm^s\mV$ is a finitely generated module over the field $\A$, and is therefore a finite dimensional $\A$-vector space.
    Let $\pi$ be any $\A$-linear map from $\frm^{s-1} \mV/\frm^s\mV$ to $\A$, then
    \[
    \sum_{i = 1}^k \ox^{p^{\ell n} \ba_i + \bb_i} \cdot \pi \left(v_i + \sum_{j = k+1}^{K'} c_{j,i}^{p^{\ell (t-1)}} \cdot \ox^{(p^{\ell (t-1)} - 1) (\bb_i - \bb_j)} v_j + \frm^s\mV \right) = 0
    \]
    for all $n \geq \max\{m, 2t\}$.
    Recall that for all $n$, the sequences
    \[
    \bx_i^{(\geq n)} = \left(\ox^{p^{\ell n} \ba_i + \bb_i}, \ox^{p^{\ell (n+1)} \ba_i + \bb_i}, \ldots\right) \in \A^{\N}
    \]
    for $i = 1, 2, \ldots, k$, are $\A$-linearly independent.
    Therefore we must have
    \[
    \pi \left(v_i + \sum_{j = k+1}^{K'} c_{j,i}^{p^{\ell (t-1)}} \cdot \ox^{(p^{\ell (t-1)} - 1) (\bb_i - \bb_j)} v_j + \frm^s\mV \right) = 0.
    \]
    for $i = 1, 2, \ldots, k$. 
    Since this is true for all linear maps $\pi \colon \frm^{s-1} \mV/\frm^s\mV \rightarrow \A$, we have
    \[
    v_i + \sum_{j = k+1}^{K'} c_{j,i}^{p^{\ell (t-1)}} \cdot \ox^{(p^{\ell (t-1)} - 1) (\bb_i - \bb_j)} v_j + \frm^s\mV = 0
    \]
    for $i = 1, 2, \ldots, k$.
    Denote
    \[
    w_i \coloneqq v_i + \sum_{j = k+1}^{K'} c_{j,i}^{p^{\ell (t-1)}} \cdot \ox^{(p^{\ell (t-1)} - 1) (\bb_i - \bb_j)} v_j
    \]
    for $i = 1, \ldots, k$, then $w_i \in \frm^s\mV$ for all $i$.
    Equation~\eqref{eq:sumyv3} can be rewritten as
    \begin{multline}\label{eq:sumyv4}
    \sum_{i = 1}^k \oA^{p^{\ell n} \ba_i + \bb_i} w_i \\
        + p \sum_{j = k+1}^{K'} \sum_{i \in \mJ} u_{i} \tc_{j,1}^{p^{\ell (t-1)}d_{i1}} \cdots \tc_{j,k}^{p^{\ell (t-1)} d_{ik}} \cdot \oA^{p^{\ell (n-t)}(\ba_1 d_{i1} + \cdots + \ba_k d_{ik}) + p^{\ell (t-1)} (\bb_1 d_{i1} + \cdots + \bb_k d_{ik})} \cdot \oA^{(1 - p^{\ell (2t-1)}) \bb_j} \cdot v_j \\
        + \sum_{i = K'+1}^K \oA^{p^{\ell n} \ba_i + \bb_i} v_i = 0.
    \end{multline}
    Note that $w_1, \ldots, w_k, pv_{k+1}, \ldots, pv_{K'}, v_{K'+1}, \ldots, v_{K}$ are now all in $\frm^s \mV$.
    The left hand side of Equation~\eqref{eq:sumyv4} is a sum of terms of the form
    $
        \oA^{p^{\ell (n-t)} \ba + \bb} v
    $,
    for $v \in \frm^s \mV$ and $\ba, \bb \in \Q^{N}$ whose denominators are not divisible by $p$.
    Indeed,
    \begin{enumerate}[nosep, label=(\roman*)]
    \item 
    the terms $\oA^{p^{\ell n} \ba_i + \bb_i} w_i, i = 1, \ldots, k$, can be written in the form $\oA^{p^{\ell (n-t)} \ba + \bb} v$, with
    \[
    \ba \coloneqq (p^{\ell n} - p^{\ell (n-t)}) \ba_i \;, \quad \bb \coloneqq \bb_i \;, \quad v \coloneqq w_i.
    \]
    Similarly, the terms $\oA^{p^{\ell n} \ba_i + \bb_i} v_i, i = K'+1, \ldots, K$, can be written in the form $\oA^{p^{\ell (n-t)} \ba + \bb} v$.
    \item
    The terms
    \[
    p \left( u_{i} \tc_{j,1}^{p^{\ell (t-1)}d_{i1}} \cdots \tc_{j,k}^{p^{\ell (t-1)} d_{ik}} \cdot \oA^{p^{\ell (n-t)}(\ba_1 d_{i1} + \cdots + \ba_k d_{ik}) + p^{\ell (t-1)} (\bb_1 d_{i1} + \cdots + \bb_k d_{ik})} \cdot \oA^{(1 - p^{\ell (2t-1)}) \bb_j} \right) \cdot v_j,
    \]
    for $j = k+1, \ldots, K'; i \in \mJ$,
    can be written in the form $\oA^{p^{\ell (n-t)} \ba + \bb} v$ with
    \[
    \ba \coloneqq \ba_1 d_{i1} + \cdots + \ba_k d_{ik} \;, \quad \bb \coloneqq p^{\ell (t-1)} (\bb_1 d_{i1} + \cdots + \bb_k d_{ik}) + (1 - p^{\ell (2t-1)}) \bb_j,
    \]
    and
    \[
    v \coloneqq p \cdot u_{i} \tc_{j,1}^{p^{\ell (t-1)}d_{i1}} \cdots \tc_{j,k}^{p^{\ell (t-1)} d_{ik}} \cdot v_j.
    \]
    \end{enumerate}
    
    Since Equation~\eqref{eq:sumyv4} can be written as a sum of terms $\oA^{p^{\ell (n-t)} \ba + \bb} v$, $v \in \frm^s \mV$, and it holds for all $n - t \geq \max\{m, 2t\} - t$, we can apply the induction hypothesis on $s$, and conclude that Equation~\eqref{eq:sumyv4} holds for all $n - t \geq (t+1-s)t$.
    Since Equation~\eqref{eq:sumyv4} is equivalent to $\sum_{i = 1}^K \oA^{p^{\ell n} \ba_i + \bb_i} v_i = 0$ for $n \geq 2t$ (which is true whenever $n - t \geq (t+1-s)t$), we conclude that 
    \[
    \sum_{i = 1}^K \oA^{p^{\ell n} \ba_i + \bb_i} v_i = 0
    \]
    holds for all $n \geq t+ (t+1-s)t = \big(t+1 - (s-1)\big) t$.
    This finishes the induction step on $s$.
\end{proof}

\end{document}
